\section{Hai xung đột mà một CEC trần tạo ra}
\label{sec:ch04-hai-xung-dot}

\index{xung đột trọng số vào!phát biểu}%
Một CEC (\tn{vòng quay lỗi không đổi}{constant error carousel}) đứng một mình
thì giữ được thông tin, nhưng nó phải nối vào phần còn lại của mạng mới có tác
dụng. Bài chỉ ra rằng nối vào là hỏng, theo hai đường đối xứng nhau.

Đường thứ nhất là \tn{xung đột trọng số vào}{input weight conflict}. Lấy một
trọng số vào \(w_{ji}\) duy nhất. Trọng số ấy phải làm hai việc: bật unit
\(j\) lên khi input đáng lưu, và để yên unit \(j\) khi input không đáng lưu.
Bài viết thẳng rằng các tín hiệu cập nhật \enquote{will attempt to make
\(w_{ji}\) participate in (1) storing the input (by switching on \(j\)) and (2)
protecting the input (by preventing \(j\) from being switched off by irrelevant
later inputs)}~\autocite{hochreiter1997lstm}. Hai việc, một con số.

\index{xung đột trọng số ra!phát biểu}%
Đường thứ hai là \tn{xung đột trọng số ra}{output weight conflict}, và nó đối
xứng: một trọng số ra \(w_{kj}\) phải vừa cho unit \(k\) đọc nội dung của \(j\)
lúc nội dung ấy có ích, vừa che \(k\) khỏi \(j\) lúc nó chỉ là nhiễu. Bài nói
thêm rằng hai xung đột này không phải đặc sản của \tn{phụ thuộc dài hạn}{long-term dependencies}:
chúng có ở mọi độ dài, chỉ nặng dần lên khi độ dài tăng, vì thông tin đã lưu
phải được bảo vệ lâu hơn và số output đã đúng cần được bảo vệ nhiều hơn.

\subsection*{Xung đột thứ nhất giải được bằng tay, nên tôi giải}

\index{xung đột trọng số vào!dạng đóng}%
Lập luận trên là lập luận của bài và nó thuyết phục, nhưng nó không cho biết
cái giá là bao nhiêu. Với một CEC trần thì cái giá viết ra được bằng giấy bút,
vì unit tuyến tính làm cả bài toán co về một phép bình phương tối thiểu một
biến.

Dựng như sau. Chuỗi dài \(T\), cell đọc mọi bước qua một trọng số vào duy nhất,
\(c_t = c_{t-1} + w x_t\), và cuối chuỗi đọc \(c_T\) rồi so với target là
\(x_1\), tức giá trị ở bước đầu tiên. Đây là bài toán \enquote{nhớ một con số
qua \(T\) bước} ở dạng trần trụi nhất. Vì cell tuyến tính,

\begin{equation}
  c_T = w \sum_{t=1}^{T} x_t = w S.
  \label{eq:ch04-cec-tran}
\end{equation}

Với \(x_t\) độc lập, kỳ vọng không, phương sai \(v\), nghiệm bình phương tối
thiểu viết ra ngay được. Sai số bình phương trung bình dưới đây viết tắt là
MSE:

\begin{equation}
  w^{*} = \frac{\E[S x_1]}{\E[S^2]} = \frac{v}{T v} = \frac{1}{T},
  \qquad
  \mathrm{MSE}(w^{*}) = v\left(1 - \frac{1}{T}\right).
  \label{eq:ch04-nghiem-toi-uu}
\end{equation}

Trọng số tốt nhất là \(1/T\), và MSE tại đó là \(v(1 - 1/T)\). Chia cho phương
sai của chính target thì phần mà cell giải thích được đúng bằng \(1/T\).

\begin{measured}{CEC trần giải thích được \(1/T\) phương sai của target}
20000 chuỗi mỗi dòng, \(x_t\) chuẩn tắc:

\begin{minted}{text}
    T  w* measured  w* = 1/T    MSE at w*   var(target)  explained
   10  0.097555    0.100000    0.888975    0.984882     0.097379
   50  0.020975    0.020000    0.975298    0.997198     0.021962
  100  0.009035    0.010000    0.994984    1.003132     0.008123
\end{minted}

Cột \code{w* measured} bám sát \(1/T\) và cột cuối bám sát \(1/T\), trong sai
số Monte Carlo ở 20000 mẫu. Ở \(T = 100\), cell giải thích 0.008123 phương sai
của target, tức dưới một phần trăm.

Đo bằng \code{experiments/ch04\_conflict.py} tại tag \code{ch04}.
\end{measured}

Con số 0.008123 ấy \emph{là} cực tiểu, và đó là chỗ đáng dừng lại. Không phải
một điểm mà optimizer kẹt lại, không phải hệ quả của learning rate tồi, không
phải chuyện huấn luyện chưa đủ lâu. Nghiệm tốt nhất tồn tại, viết ra được, và
nó vô dụng. Một kiến trúc mà lời giải tối ưu cũng không dùng được thì không sửa
được bằng thuật toán huấn luyện, và đó chính là lý do chương này tồn tại tách
khỏi chương~\ref{ch:phan-tich}.

Xung đột thứ hai tôi không đo. Lập luận cho nó là lập luận của bài, đối xứng
với lập luận vừa rồi, và tôi để nguyên nó ở dạng bài viết thay vì dựng một số
đo mà thiết kế thí nghiệm sẽ quyết định phần lớn kết quả.

\subsection*{Cái thiếu là một chỗ để đặt chữ \enquote{khi nào}}

Nhìn lại~\eqref{eq:ch04-cec-tran} thì chỗ hỏng lộ ra. Trọng số \(w\) là một
hằng số theo thời gian, còn câu trả lời đúng thì không: nó phải là 1 tại bước
1 và 0 ở mọi bước khác. Không có giá trị hằng nào của \(w\) biểu diễn được một
hàm theo \(t\), nên bình phương tối thiểu trả về trung bình của thứ nó cần, tức
\(1/T\).

Vậy cái thiếu không phải một trọng số tốt hơn mà là một thừa số \emph{thay đổi
theo thời gian}, nhân vào input, do chính mạng tính ra từ ngữ cảnh. Đặt một
thừa số như thế vào và cùng bài toán đó thành tầm thường:

\begin{minted}{text}
    T  MSE at w=1
   10  0.000000e+00
   50  0.000000e+00
  100  0.000000e+00
\end{minted}

Sai số bằng không, ở cả ba độ dài, vì \(c_T = w x_1\) và \(w = 1\) là nghiệm
đúng. Cần nói rõ tôi đã làm gì để có bảng này: tôi \emph{đặt} thừa số ấy bằng
tay, mở ở bước 1 và đóng sau đó, chứ không huấn luyện ra nó. Bảng này không
chứng minh mạng học được lời giải; nó chứng minh kiến trúc biểu diễn được lời
giải, và bảng trước đã chứng minh kiến trúc kia thì không.

Thừa số nhân theo thời gian ấy là cái mà bài gọi là cổng.
